\documentclass[a4paper,14pt]{extarticle}

\usepackage[T2A]{fontenc}
\usepackage[utf8]{inputenc}
\usepackage[russian]{babel}

% Геометрия страницы (по ГОСТ: левое 30 мм, правое 15 мм, верхнее 20 мм, нижнее 20 мм)
\usepackage{geometry}
\geometry{left=3cm, right=1.5cm, top=2cm, bottom=2cm}

\usepackage{amsmath, amsfonts, amssymb}
\usepackage{graphicx}
\usepackage{booktabs}
\usepackage{setspace}
\usepackage{titlesec}
\usepackage{indentfirst}

% Каждая секция (\section) всегда начинается с новой страницы
\newcommand{\sectionbreak}{\clearpage}

% Оформление заголовков секций
\titleformat{\section}{\singlespacing\normalfont\Large\bfseries}{\thesection}{1em}{}
\titleformat{\subsection}{\singlespacing\normalfont\large\bfseries}{\thesubsection}{1em}{}

% Настройка гиперссылок
\usepackage{hyperref}
\hypersetup{
    colorlinks=true,
    linkcolor=black,
    filecolor=blue,      
    urlcolor=blue,
    pdfpagelabels=false
}

% --- Настройка стилей ---
\onehalfspacing % Полуторный интервал
\setlength{\parindent}{1.25cm} % Абзацный отступ
\emergencystretch=3em % Защита от вылезания слов за границы полей

\begin{document}
\sloppy

% --- Титульный лист ---
\begin{titlepage}
    \begin{center}
        Министерство науки и высшего образования РФ\\
        \vspace{0.4cm}
        Федеральное государственное автономное образовательное учреждение\\
        высшего образования «Казанский (Приволжский) федеральный университет»\\
        \vspace{0.4cm}
        ИНСТИТУТ ВЫЧИСЛИТЕЛЬНОЙ МАТЕМАТИКИ 
        И ИНФОРМАЦИОННЫХ ТЕХНОЛОГИЙ\\
        КАФЕДРА ПРИКЛАДНОЙ МАТЕМАТИКИ
        И ИСКУССТВЕННОГО ИНТЕЛЛЕКТА\\
        Направление: 01.03.04 – Прикладная математика\\
        Профиль: Прикладная математика

        
        \vspace{3.5cm}
        \textbf{\Large ОТЧЕТ ПО ПРОЕКТУ}\\
        \vspace{0.5cm}
        \textbf{\large «Многослойный персептрон»}\\
        \vspace{0.8cm}
        По курсу «Анализ на Python»
    \end{center}

    \vspace{2.5cm}
    \begin{flushright}
        \textbf{Выполнил:}\\
        Студент группы 09-321\\
        Суюндуков Г.Э.\\
        \vspace{0.8cm}
        \textbf{Проверил:}\\
        Старший преподаватель\\
        Гиниятова Д.Х.\\
    \end{flushright}

    \vfill
    \begin{center}
        Казань, 2026
    \end{center}
\end{titlepage}

% --- Содержание ---
\clearpage
\pagenumbering{arabic}
\setcounter{page}{2}
\tableofcontents
\clearpage

% --- Введение ---
\section{Введение}
Целью данного проекта является реализация искусственной 
нейронной сети прямого распространения (многослойного персептрона) 
для решения задачи бинарной классификации. 
Все алгоритмы, включая прямое и обратное распространение ошибки,
а также оптимизаторы SGD и Adam, реализованы исключительно с использованием библиотеки линейной алгебры \texttt{NumPy}.
В качестве набора данных используется датасет Breast Cancer Wisconsin. Набор содержит 569 образцов и 30 числовых признаков, описывающих характеристики ядер клеток. Целевая переменная (диагноз) имеет два класса: \textbf{M} (злокачественная) и \textbf{B} (доброкачественная). 

Архитектура проекта состоит из трех модулей:
\begin{enumerate}
    \item \texttt{split.py} — разведочный анализ (EDA), предобработка и разделение данных;
    \item \texttt{train.py} — реализация и обучение нейронной сети;
    \item \texttt{predict.py} — прогнозирование на тестовой выборке и оценка качества модели.
\end{enumerate}

% --- Предобработка ---
\section{Разведочный анализ и предобработка данных}
\subsection{Разведочный анализ (EDA)}
Анализ распределения классов показал наличие дисбаланса (рис. \ref{fig:class_dist}): 357 образцов 
(62.7\%) относятся к классу доброкачественных опухолей, и 212 образцов 
(37.3\%) — к злокачественным.

\begin{figure}[h!]
    \centering
    \includegraphics[width=0.8\textwidth]{plot_class_distribution.png}
    \caption{Распределение классов доброкачественных (B) и злокачественных (M) опухолей в датасете.}
    \label{fig:class_dist}
\end{figure}

Для визуального анализа признаков были 
построены гистограммы и диаграммы размаха, позволившие выявить признаки, обладающие наибольшей разделяющей способностью (например, \texttt{concavity\_mean}). Также была вычислена матрица корреляций Пирсона $R$:
\begin{equation}
    R = \frac{\tilde{X}^T \tilde{X}}{n}
\end{equation}
где $\tilde{X}$ — центрированная и нормированная матрица признаков.

\subsection{Стратифицированное разделение и устранение пропусков}
Данные были разделены на обучающую (80\%) и валидационную (20\%) выборки с сохранением пропорций классов (стратификацией). Для предотвращения \textit{утечки данных (data leakage)} разделение производилось до этапа нормализации и заполнения пропусков.

Заполнение возможных пропущенных значений осуществлялось не средним значением по всему датасету, а средним значением по соответствующему классу, вычисленным \textbf{только на обучающей выборке}. Это позволяет сохранить распределения признаков для каждого класса и не допустить смешивания характеристик доброкачественных и злокачественных клеток.

\subsection{Нормализация признаков}
Признаки имеют различные масштабы (например, \texttt{area\_mean} достигает 2500, в то время как \texttt{smoothness\_mean} находится в диапазоне 0.05-0.16). Для обеспечения стабильной и быстрой сходимости градиентного спуска применялась Z-score нормализация:
\begin{equation}
    z = \frac{x - \mu}{\sigma}
\end{equation}
Параметры $\mu$ и $\sigma$ вычислялись строго по обучающей выборке и затем применялись к валидационной. 

Целевая переменная была закодирована: M $\rightarrow$ 1, B $\rightarrow$ 0.

% --- Архитектура ---
\section{Архитектура нейронной сети и алгоритм обучения}
\subsection{Прямое распространение (Forward Pass)}
Был реализован многослойный персептрон (MLP), принимающий на вход 30 признаков. Формула вычисления активаций $l$-го слоя:
\begin{equation}
    Z^{[l]} = A^{[l-1]} W^{[l]} + b^{[l]}, \quad A^{[l]} = g(Z^{[l]})
\end{equation}
В качестве функции активации скрытых слоев используется сигмоида:
\begin{align}
    \sigma(z) &= \frac{1}{1 + e^{-z}} \\
    \sigma'(z) &= \sigma(z)(1 - \sigma(z))
\end{align}
На выходном слое (2 нейрона) применяется функция Softmax для получения вероятностного распределения по классам:
\begin{equation}
    \text{softmax}(z_k) = \frac{e^{z_k - \max z}}{\sum_j e^{z_j - \max z}}
\end{equation}
Вычитание максимума обеспечивает числовую стабильность вычислений.

Для инициализации весовых матриц использовался метод He Uniform:
\begin{equation}
    W \sim U\left(-\sqrt{\frac{6}{n_{\text{in}}}}, \sqrt{\frac{6}{n_{\text{in}}}}\right)
\end{equation}

\subsection{Функция потерь и обратное распространение}
В качестве функции потерь выбрана Categorical Cross-Entropy:
\begin{equation}
    \mathcal{L} = -\frac{1}{m} \sum_{i=1}^{m} \sum_{k=0}^{1} y_{ik} \log(\hat{p}_{ik})
\end{equation}
где $y_{ik}$ — one-hot представление истинной метки класса, $\hat{p}_{ik}$ — предсказанная Softmax-вероятность.

Комбинация Softmax и Cross-Entropy дает численно стабильное выражение для градиента на последнем слое:
\begin{equation}
    \delta^{[L]} = \frac{\partial \mathcal{L}}{\partial Z^{[L]}} = \hat{p} - y_{\text{one-hot}}
\end{equation}
Градиенты для скрытых слоев вычисляются по правилу цепочки:
\begin{equation}
    \delta^{[l-1]} = \left(\delta^{[l]} (W^{[l]})^T\right) \odot \sigma'(Z^{[l-1]})
\end{equation}
Градиенты весов и смещений:
\begin{align}
    \frac{\partial \mathcal{L}}{\partial W^{[l]}} &= \frac{1}{m}(A^{[l-1]})^T \delta^{[l]} \\
    \frac{\partial \mathcal{L}}{\partial b^{[l]}} &= \frac{1}{m}\sum \delta^{[l]}
\end{align}

\subsection{Оптимизаторы: SGD и Adam}
Для обновления весов реализовано два алгоритма.

\textbf{Мини-батч стохастический градиентный спуск (SGD)}:
\begin{equation}
    W \leftarrow W - \eta \frac{\partial \mathcal{L}}{\partial W}
\end{equation}

\textbf{Adam (Adaptive Moment Estimation)}:
\begin{align}
    m_t &= \beta_1 m_{t-1} + (1 - \beta_1) g_t \\
    v_t &= \beta_2 v_{t-1} + (1 - \beta_2) g_t^2 \\
    \hat{m}_t &= \frac{m_t}{1 - \beta_1^t} \\
    \hat{v}_t &= \frac{v_t}{1 - \beta_2^t} \\
    W &\leftarrow W - \frac{\eta}{\sqrt{\hat{v}_t} + \epsilon} \hat{m}_t
\end{align}
Алгоритм Adam адаптирует шаг обучения для каждого параметра, что значительно ускоряет сходимость.

% --- Результаты ---
\section{Оценка качества модели}
Обучение проводилось на архитектуре с двумя скрытыми слоями (64 и 32 нейрона) размером мини-батча 32. 
Основная используемая метрика для бинарной классификации — Binary Cross-Entropy (BCE). Для детального анализа использовались метрики Accuracy, Precision, Recall и F1-score.

\subsection{Сравнение оптимизаторов}
Сравнение эффективности оптимизаторов SGD ($\eta = 0.01$) и Adam ($\eta = 0.001$) показало превосходство адаптивного алгоритма:

\begin{table}[h!]
    \centering
    \begin{tabular}{l c c c}
        \toprule
        \textbf{Оптимизатор} & \textbf{Эпохи} & \textbf{Accuracy (Val)} & \textbf{BCE Loss (Val)} \\
        \midrule
        SGD & 50 & 93.81\% & 0.264 \\
        Adam & 100 & \textbf{98.23\%} & \textbf{0.114} \\
        \bottomrule
    \end{tabular}
    \caption{Сравнение оптимизаторов на валидационной выборке}
\end{table}

\subsection{Анализ матрицы ошибок}
Для лучшей модели (Adam) была построена матрица ошибок (Confusion Matrix) на валидационной выборке ($N=113$), представленная в таблице 2 и на рисунке \ref{fig:conf_matrix}.

\begin{table}[h!]
    \centering
    \begin{tabular}{l c c}
        \toprule
         & \textbf{Предсказан B} & \textbf{Предсказан M} \\
        \midrule
        \textbf{Истинный B} & 71 (TN) & 0 (FP) \\
        \textbf{Истинный M} & 2 (FN) & 40 (TP) \\
        \bottomrule
    \end{tabular}
    \caption{Матрица ошибок}
\end{table}

\begin{figure}[h!]
    \centering
    \includegraphics[width=0.65\textwidth]{plot_confusion_matrix.png}
    \caption{Визуализация матрицы ошибок для модели, обученной оптимизатором Adam.}
    \label{fig:conf_matrix}
\end{figure}

Результаты показывают отсутствие ложных срабатываний (Precision для класса M равен 1.0). Из 42 злокачественных опухолей модель успешно выявила 40 (Recall = 0.952), F1-мера для целевого класса составила 0.975. Отсутствие переобучения подтверждается тем, что кривые обучения (Loss и Accuracy) для обучающей и валидационной выборок сходятся на графиках синхронно (рис. \ref{fig:learning_curves}).

\begin{figure}[h!]
    \centering
    \includegraphics[width=1\textwidth]{plot_learning_curves.png}
    \caption{Кривые обучения функции потерь (слева) и точности (справа) для оптимизатора Adam. Сходимость тренировочных и валидационных кривых указывает на отсутствие переобучения.}
    \label{fig:learning_curves}
\end{figure}

\section{Заключение}
В ходе выполнения проекта была с нуля разработана полностью 
функциональная библиотека для создания нейронных сетей, 
поддерживающая абстракции слоев, функций активации и оптимизаторов. 
Модель многослойного персептрона успешно решила задачу классификации
 клеток молочной железы, достигнув точности в 98.23\% на отложенной 
 выборке. Использование современных подходов (He Uniform, 
 Softmax+Cross-Entropy, нормализация без утечки данных, Adam) 
 доказало свою эффективность.

\end{document}
